\chapter{Lernziele}
\begin{todolist}
    \item Ich kann erklären, was ein Bit und ein Byte ist, und gängige Dateigrösseneinheiten (KB, MB, GB, TB) angeben.
    \item Ich kann die Wahrheitstabellen der logischen Grundgatter NOT, AND und OR aufstellen.
    \item Ich kann die Symbole der Gatter NOT, AND, OR und NAND zeichnen und beschriften.
    \item Ich kann das NAND-Gatter als Wahrheitstabelle angeben und erklären, warum es funktional vollständig ist.
    \item Ich kann den Halbaddierer aus AND- und XOR-Gattern zusammensetzen und seine Funktion (Summe und Übertrag) erklären.
    \item Ich kann die vier Hauptkomponenten der Von-Neumann-Architektur (CPU, RAM, Massenspeicher, Ein-/Ausgabe) benennen und ihre Aufgaben beschreiben.
    \item Ich kann die Bestandteile der CPU (ALU, Steuerwerk, Register) erläutern.
    \item Ich kann den Fetch-Decode-Execute-Zyklus mit eigenen Worten beschreiben.
    \item Ich kann den Unterschied zwischen RAM und Massenspeicher (HDD, SSD) bezüglich Geschwindigkeit, Flüchtigkeit und Grösse erklären.
    \item Ich kann das Schichtenmodell (Software-Stack) mit den Ebenen Hardware, Treiber, Betriebssystem, API und Anwendung skizzieren und erklären.
    \item Ich kann die vier zentralen Aufgaben eines Betriebssystems (Prozessverwaltung, Speicherverwaltung, Dateisystem, Geräteverwaltung) beschreiben.
    \item Ich kann das Prinzip des Multitasking (Zeitscheiben) erläutern.
    \item Ich kann an einem konkreten Beispiel (z.\,B. Tastendruck) das Zusammenspiel von Hardware, Betriebssystem und Anwendung erklären.
\end{todolist}
